Large language models (LLMs)~\citep{brown2020language, gpt-j, zhang2022opt, smith2022using,  chowdhery2022palm, hoffmann2022training, workshop2023bloom, openai2023gpt4, anil2023palm, touvron2023llama, touvron2023llama2, penedo2023refinedweb} have significantly advanced the state-of-the-art on a wide range of language understanding tasks, leading to widespread deployment and adoption in applications~\citep{araci2019finbert, huang2020clinicalbert, biomedlm, wu2023bloomberggpt, driess2023palme, huang2023instruct2act}. In particular, the emerging capabilities of LLMs such as complex reasoning~\citep{wei2023chainofthought, wang2023selfconsistency, kıcıman2023causal} have made them the subject of research in recent years. Among these, zero-shot reasoning~\citep{kojima2023large,wan2023better} has drawn substantial public interest and achieved promising results in specific task domains. However, the reasoning ability of LLMs on general tasks remains unclear.

\begin{figure*}[ht!]
  \centering
  \includegraphics[width=\textwidth]{figures/fig1-once-per-5-31-final.drawio.png}
  \caption{{Summary of our approach and results. Top: \zeroshotstart\ \method\ generalizes the zero-shot reasoning abilities of large language models to a wide set of language understanding tasks including generation, classification, and reasoning. Our agent runs only once per task, producing just one set of task-specific instructions that are used with all $n$ instances of the task to instruct the reasoning process of large language models. These same instructions are also shared across all $m$ models.
  As an example, results from instance 1 and model $m$ are shown.
  Both the agent instructions and task-specific reasoning process are highlighted. 
  Bottom: Performance of \zeroshotmid\ \method\ compared with standard zero-shot and zero-shot chain of thought (CoT). \zeroshotstart\ \method\ improves the performance of three large language models substantially on the 29 datasets we evaluate.
  }}
  \label{fig:full-pipeline}
\end{figure*}

In this paper, we improve the zero-shot reasoning abilities of LLMs on general language understanding tasks. To solve a task, we build an agent to instruct the reasoning process of LLMs for the task (Figure~\ref{fig:full-pipeline}). For each task, our autonomous agent generates a unique set of task-specific instructions. These instructions, produced only once per task, are used to guide the reasoning of a variety of LLMs across all task instances. Our agent is built upon a larger LLM teaching the reasoning process to smaller LLMs. This design follows the general knowledge distillation setup involving the teacher-student paradigm. Intuitively, task-specific instructions help better align the chain of thought reasoning process of LLMs with each task. We refer to this approach as zero-shot agent instructed reasoning (\method). The basic idea of our approach is motivated by two lines of work. First, the development of language agents~\citep{yao2023react, shinn2023reflexion, park2023generative, wang2023survey, xi2023rise} to automatically complete a task. Instead of completing the task, our agent produces instructions on how to complete the task. We enable this by adapting an existing agent to access a wide variety of task-relevant knowledge on the web, given basic task information such as the dataset name and several input-only examples. As a result, the agent synthesizes high-quality step-by-step instructions for tasks, verified by web resources. We follow the recent design of language agents that plan a process, though our process runs only once per dataset instead of per dataset instance. Second, zero-shot chain of thought (CoT) reasoning of LLMs has obtained promising results on tasks such as arithmetic reasoning~\citep{kojima2023large, wang2023planandsolve}. Standard zero-shot learning prompts an LLM to directly output predictions without task examples. In contrast, CoT decomposes a task into intermediate steps where solving each will lead to the final output. We further align the CoT reasoning steps with a particular task by prompting with task-specific agent instructions. The design of \zeroshotmid\ \method\ is important: We generalize the zero-shot reasoning abilities of LLMs to more tasks by combining task-specific instructions from a language agent and task-specific reasoning of LLMs.

We empirically evaluate the zero-shot reasoning abilities of LLMs on a wide set of language understanding tasks across \NumDatasets\ datasets (including \NumSubsets\ subsets), spanning generation, classification, and reasoning. \zeroshotstart\ \method\ obtains state-of-the-art performance on 20 datasets. We conduct our evaluation on three state-of-the-art LLMs, namely, Vicuna \citep{vicuna2023}, Llama-2-chat \citep{touvron2023llama2}, and GPT-3.5 Turbo \citep{openaichatgpt}. We show that \zeroshotmid\ \method\ boosts the performance of these models by \ThreemodelsOverallAvgwinOverZs\ on average. When compared to zero-shot CoT, the overall performance improvement is significant (\ThreemodelsOverallAvgwinOverZscot), and in particular, the improvement in reasoning is substantial with an average increase of \ThreemodelsReasoningAvgwinOverZscot, leading to the best performance on 10 out of 12 reasoning tasks. Notably, Llama-2-70b-chat with \zeroshotmid\ \method\ outperforms standard zero-shot GPT-3.5 Turbo by an average of \LlamaTwoSeventyChatAvgwinOverChatgptZs. We hope the results help foster future research on further unlocking the zero-shot reasoning abilities of large foundation models and exploring the broader usage of agents.